\bookchapter{附錄C：極簡術語“黑話”對照表}

\begin{center}
\begin{tabularx}{\linewidth}{|X|X|X|}
\hline
\textbf{你可能聽到的詞} & \textbf{其實就是指……} & \textbf{一句話解釋} \tabularnewline
\hline
SoC / 晶片 & 手機的大腦和心臟 & 決定手機快不快、遊戲卡不卡的核心部件。 \tabularnewline
RAM (運存) & 手機的"短期工作桌" & 桌越大，能同時開的App就越多，切換越流暢。 \tabularnewline
ROM (儲存) & 手機的"大倉庫" & 存照片、影片、App的地方。128G、256G說的就是它。 \tabularnewline
OLED / LCD & 螢幕的兩種材質 & OLED色彩更豔、更省電；LCD更護眼、成本低。 \tabularnewline
NFC & 近距離無線通訊 & 讓你手機能變身公交卡、門禁卡、銀行卡的"貼貼"功能。 \tabularnewline
OTA更新 & 無線系統升級 & 在【設定】裡點的那個"系統更新"，就是它。 \tabularnewline
Root & 安卓超級管理員許可權 & 獲得後可以修改系統檔案，但也有風險。 \tabularnewline
BL鎖 & Bootloader鎖 & 廠商加的安全鎖，解鎖後才能刷機。 \tabularnewline
NPU & 神經網路處理器 & 專門處理AI任務的晶片，比CPU快幾十倍、功耗更低。 \tabularnewline
端側AI & 手機本地執行的AI & AI在手機上直接執行，不需要聯網，速度快、隱私好。 \tabularnewline
Qi2 & 新一代無線充電標準 & 帶磁吸定位的無線充電，類似MagSafe，充電更方便。 \tabularnewline
eSIM & 嵌入式SIM卡 & 沒有實體卡的SIM，直接寫入手機，換運營商不用換卡。 \tabularnewline
LTPO & 自適應重新整理率技術 & 螢幕重新整理率根據內容自動調節（1-120Hz），流暢又省電。 \tabularnewline
\hline
\end{tabularx}
\end{center}
